\chapter{Implications and Limitations}
\label{chap:discussion}

\section{Implications for using Expanded Welch}

Expanded Welch improves false-positive control in several small-sample
settings. For uniform \(4\times4\) populations at \(n=20\), it lowers
the false-positive rate from 0.0813 to 0.0470, with every calculation valid.
For \(8\times8\) different-skew populations at \(n=100\), it lowers the
rate from 0.0706 to 0.0518, again with full validity.

The power comparison uses larger samples, at \(n=1{,}000\) and 10,000,
so it does not establish the power cost in those smaller-sample settings.
Across the power comparisons, Expanded Welch rejects less often and offers
no general advantage over Normal Wald. It also lowers false-positive rates
that are already conservative and does not correct the unequal biases in
the MI difference. Normal Wald therefore remains the simpler
baseline. Valid-result frequency remains important, because a lower
rejection rate can otherwise reflect failed calculations.

For an observed feature--label comparison, report the estimated MI
difference, its standard error and the p-value together. The difference
gives the estimated size and direction of the change, whereas the standard error
gives the sampling uncertainty assigned by the method. The sparse-table
results show why the direction matters: unequal estimation biases can
shift or reverse the estimated difference, as
Equations~\eqref{eq:focused-null-bias-shift} and
\eqref{eq:focused-detection-bias-shift} illustrate. If Expanded
Welch is used, report its estimated degrees of freedom as well, since
they determine the rejection cutoff for that comparison.

Before applying either test, assess its performance on populations that
match the application's category definitions, sample sizes and margin
patterns. Include plausible population MI values and distinct equal-MI
tables, then compare false positives, power and usable-result frequency.
Because neither method is reliable in every studied setting, this
assessment is particularly important for sparse data. Choose the method using
that evidence before examining the application's observed p-values.

\begin{samepage}
\section{Limits of the evidence}

The rejection rates and Monte Carlo intervals apply to the constructed
population pairs. The intervals quantify uncertainty from the finite number
of simulated samples; they do not describe how performance varies across
untested populations.
\par
\end{samepage}

The primary populations use a copying mixture with identical row and
column margins within each population, allowing association to vary
without changing category frequencies. The four-cell check
supports the representative \(4\times4\) method comparison and
reproduces the main \(16\times16\) estimation failures at its matched
different-skew settings, although power changes numerically with the construction.
The same-skew checks instead give conservative sparse null rates and low
power under both constructions. With equal MI and equal sample sizes,
matching category frequencies gives identical population tables within
each construction, so the expected MI estimation biases cancel. This does
not ensure accurate estimation of uncertainty.
Copying favours matching categories in both
populations, whereas the four-cell transfer concentrates association
in the same chosen block. The different-skew regime changes category
frequencies while keeping those association locations aligned. Its
results therefore demonstrate that distinct equal-MI tables can be
difficult to compare, but do not describe every way two tables with
equal MI can differ. A broader assessment would vary association
locations, population MI values and margins, including unequal
row and column margins within a population.

Because equal sample sizes make the two leading bias corrections cancel
in \(\Deltahat\), the remaining difference bias in the main studies
is exactly the bias of the raw plug-in MI difference. This conclusion
concerns the fixed-alphabet correction used here. With unequal sample sizes, the
corrections no longer cancel, and the two contributions to the standard
error and degrees of freedom can become less balanced. The derivation
allows unequal sample sizes. The supporting fixed-total check assesses
both allocation directions in four \(4\times4\) pairs: a null and an
alternative under each skew regime.
There, imbalance changes bias, uncertainty and power, while the
larger Expanded Welch cutoff lowers already-conservative null rates.
The same-skew null shows that unequal sample sizes can leave unequal
residual biases even when the two population tables are identical.
This extends the comparison beyond equal samples, but a wider allocation
study would be needed for other table sizes and population structures.

The sampling assumptions also limit transfer to other applications.
Paired observations would require a covariance term in the variance of
the difference, whereas the present samples are independent. Continuous
MI estimators have different estimation errors from the categorical MI
estimator used here. Similarly, increasing the number of categories with
sample size changes the fixed-alphabet setting of the derivation. The empirical
recommendation concerns the independent categorical comparisons studied
here, at a two-sided 5\% threshold. The online-shopping application
illustrates interpretation on observed data, but its unknown population
MI values do not allow a direct measurement of false-positive rates or power.

Exact independence needs a different mathematical approximation because
the first-order MI variance vanishes there. The population-MI
study also shows that positive population MI is not enough to make that
approximation accurate at a particular sample size: small positive MI
can remain close to the problematic region. This explains why improving
the rejection cutoff and improving MI's sampling approximation near independence
are separate tasks.
Appendix~\ref{sec:independence-boundary} derives the second-order
approximation needed there.
